% Source-faithful assembly: Mumford, Abelian Varieties, Chapter IV,
% Horizon transcription pages 0197--0206.

\section{Riemann forms}
\label{mumford-IV20-source}
\dcref{M:ch4:IV.20}

\begin{proposition}[Tate pairing and Riemann form]
\label{mumford-prop-tate-riemann-pairing}
\dcref{M:ch4:IV.20}
\uses{mumford-def-weil-pairing-ladic}
There is a natural non-degenerate $\mathbb Z_l$-bilinear pairing
\[
e_l:T_l(X)\times T_l(\widehat X)\longrightarrow M_l.
\]
If $f:X\to Y$ has dual $\widehat f$, then
\[
e_l(T_l(f)\xi,\eta)=e_l(\xi,T_l(\widehat f)\eta).
\]
For a line bundle $L$ on $X$, define
\[
E^L(x,y)=e_l(x,T_l(\phi_L)(y)).
\]
\source{mumford-abelian-varieties:page-0197}
\end{proposition}

\begin{theorem}[Skew-symmetry of the Riemann form]
\label{mumford-thm-riemann-form-skew}
\dcref{M:ch4:IV.20}
\uses{mumford-prop-tate-riemann-pairing}
The form $E^L$ is skew-symmetric for every line bundle $L$.
Consequently $L\mapsto E^L$ injects
\[
\operatorname{NS}(X)\longrightarrow
\{\text{alternating forms }T_l(X)\times T_l(X)\to M_l\}.
\]
\source{mumford-abelian-varieties:page-0197}
\source{mumford-abelian-varieties:page-0198}
\end{theorem}
\begin{proof}
It suffices to show $e_l(a,\phi_L(a))=1$ for $a\in X_l$. If $a\in X_n$
and $(g)=n_X^{-1}(T_aD-D)$ for a divisor $D$ representing $L$, choose
$b$ with $nb=a$. Then $h(x)=\prod_{i=0}^{n-1}g(x+ib)$ is constant, and
\[
1=\frac{h(x+b)}{h(x)}=\frac{g(x+a)}{g(x)}.
\]
Thus the pairing is trivial on the diagonal. If $E^L=0$, then
$T_l(\phi_L)=0$, and faithfulness of the Tate representation gives
$\phi_L=0$, proving injectivity modulo $\operatorname{Pic}^0$.
\end{proof}

\begin{proposition}[Pullback of Riemann forms]
\label{mumford-prop-riemann-form-pullback}
\dcref{M:ch4:IV.20}
\uses{mumford-prop-tate-riemann-pairing}
For a homomorphism $f:X\to Y$ and a line bundle $L$ on $Y$,
\[
E^{f^*L}(x,y)=E^L(T_l(f)x,T_l(f)y).
\]
\source{mumford-abelian-varieties:page-0198}
\end{proposition}
\begin{proof}
Using $\phi_{f^*L}=\widehat f\phi_Lf$ and adjointness,
\[
\begin{aligned}
E^{f^*L}(x,y)&=e_l(x,\phi_{f^*L}(y))\\
&=e_l(x,T_l(\widehat f\phi_Lf)(y))
=e_l(T_l(f)x,\phi_L(T_l(f)y)).
\end{aligned}
\]
\end{proof}

\begin{proposition}[The Poincare Riemann form]
\label{mumford-prop-poincare-riemann-form}
\dcref{M:ch4:IV.20}
\uses{mumford-thm-riemann-form-skew,mumford-thm-poincare-universal-mapping}
Under $T_l(X\times\widehat X)=T_l(X)\times T_l(\widehat X)$,
\[
E^P((x,\widehat x),(y,\widehat y))
=e_l(x,\widehat y)-e_l(y,\widehat x).
\]
\source{mumford-abelian-varieties:page-0198}
\source{mumford-abelian-varieties:page-0199}
\end{proposition}
\begin{proof}
By skew-symmetry and linearity it suffices to consider the two factors
separately. The first two restrictions of $P$ are trivial. The canonical
identification $\widehat{X\times Y}\simeq\widehat X\times\widehat Y$ gives
$\phi_P(x,\widehat x)=(\widehat x,\phi(x))$, and hence the cross term is
$e_l(x,\widehat y)$; skew-symmetry supplies the other cross term.
\end{proof}

\begin{theorem}[Skew forms and line bundles]
\label{mumford-thm-riemann-form-line-bundle}
\dcref{M:ch4:IV.20}
\uses{mumford-prop-poincare-riemann-form}
For a homomorphism $\phi:X\to\widehat X$, the form
$e_l(x,\phi(y))$ is skew-symmetric if and only if there is a line bundle
$L$ with $2\phi=\phi_L$.
\source{mumford-abelian-varieties:page-0199}
\end{theorem}
\begin{proof}
The forward implication follows from Theorem 1. Conversely, pull back $P$
along $(1,\phi):X\to X\times\widehat X$ to obtain $L$. Then
\[
e_l(x,\phi_L(y))=E^P((x,\phi x),(y,\phi y))
=e_l(x,\phi y)-e_l(y,\phi x)=2e_l(x,\phi y),
\]
so faithfulness gives $2\phi=\phi_L$.
\end{proof}

\begin{proposition}[Rosati involution]
\label{mumford-prop-rosati-involution}
\dcref{M:ch4:IV.20}
\uses{mumford-prop-tate-riemann-pairing}
For ample $L$ on $X$, define
\[
\phi'=\phi_L^{-1}\widehat\phi\phi_L.
\]
Then $'$ is a $\mathbb Q$-linear involution,
$(\phi\psi)'=\psi'\phi'$, and
\[
E^L(\phi x,y)=E^L(x,\phi'y).
\]
\source{mumford-abelian-varieties:page-0200}
\end{proposition}
\begin{proof}
Extend $T_l$ to $\operatorname{End}^0X$. The displayed identity follows
from the Tate adjunction:
\[
E^L(x,\phi'y)=e_l(x,\widehat\phi\phi_L(y))
=e_l(\phi x,\phi_L(y))=E^L(\phi x,y).
\]
Non-degeneracy gives all involution identities.
\end{proof}

\begin{theorem}[Neron--Severi classes and Rosati fixed points]
\label{mumford-thm-ns-rosati-fixed}
\dcref{M:ch4:IV.20}
\uses{mumford-thm-riemann-form-line-bundle,mumford-prop-rosati-involution}
Under $\operatorname{Hom}^0(X,\widehat X)\simeq\operatorname{End}^0X$ via
$\psi\mapsto\phi_L^{-1}\psi$, the space
$\mathbb Q\otimes\operatorname{NS}(X)$ is exactly the fixed subspace of
the Rosati involution.
\source{mumford-abelian-varieties:page-0201}
\end{theorem}
\begin{proof}
By the preceding skew-form theorem, $\psi$ comes from a divisor precisely
when $e_l(x,\psi y)=-e_l(y,\psi x)$. This is equivalent to
\[
E^L(x,\psi y)=E^L(\psi x,y)=E^L(x,\psi'y).
\]
Non-degeneracy and faithfulness give $\psi=\psi'$.
\end{proof}

\begin{theorem}[Wedge and intersection numbers]
\label{mumford-thm-riemann-wedge-intersection}
\dcref{M:ch4:IV.20}
\uses{mumford-thm-riemann-roch,mumford-thm-degree-tate-determinant}
There is a generator
$v\in\operatorname{Hom}_{\mathbb Z_l}(\Lambda^2T_l(X),M_l)$ such that
for divisors $D_i$ and $E_i=E^{\mathcal O_X(D_i)}$,
\[
E_1\wedge\cdots\wedge E_g=(D_1\cdots D_g)v.
\]
\source{mumford-abelian-varieties:page-0201}
\source{mumford-abelian-varieties:page-0202}
\end{theorem}
\begin{proof}
Both sides are polynomial in the Neron--Severi classes, so polarization
reduces to one bundle $L$. Riemann--Roch gives
$[E^L]^{\wedge g}=g!\chi(L)v$. In a basis,
\[
([E^L]^{\wedge g})^2=(\det E^L)(g!)^2.
\]
Since $\det E^L=\det T_l(\phi_L)$ and the degree/determinant comparison
differs only by an $l$-adic unit, the normalization of $v$ gives the result.
\end{proof}

\begin{corollary}[Degree of a fixed subalgebra]
\label{mumford-cor-rosati-subalgebra-degree}
\dcref{M:ch4:IV.20}
\uses{mumford-thm-riemann-wedge-intersection}
If $X$ is simple and $K\subset\operatorname{End}^0X$ is a $\mathbb Q$-subalgebra
fixed by Rosati, then $[K:\mathbb Q]$ divides $\dim X$.
\source{mumford-abelian-varieties:page-0202}
\source{mumford-abelian-varieties:page-0203}
\end{corollary}
\begin{proof}
The ratio $\chi(M)/\chi(L)$ is a norm form on $K$ of degree $g$, since its
square is $\deg(\phi_L^{-1}\phi_M)$. The degree of a field therefore divides
the degree of this norm form.
\end{proof}

\section{Positivity of the Rosati involution}
\label{mumford-IV21-source}
\dcref{M:ch4:IV.21}

\begin{theorem}[Positive Rosati trace]
\label{mumford-thm-rosati-positive}
\dcref{M:ch4:IV.21}
\uses{mumford-thm-riemann-wedge-intersection,mumford-prop-rosati-involution}
For an ample divisor $H$, $L=\mathcal O_X(H)$, and the Rosati involution
defined by $L$,
\[
\operatorname{Tr}(\phi\phi')
=\frac{2g}{(H^g)}(H^{g-1}\cdot\phi^*H),
\]
and $\phi\mapsto\operatorname{Tr}(\phi\phi')$ is positive definite.
\source{mumford-abelian-varieties:page-0203}
\source{mumford-abelian-varieties:page-0204}
\end{theorem}
\begin{proof}
The wedge/intersection theorem gives
\[
\frac{(E^L)^{\wedge(2g-1)}\wedge E^{\phi^*L}}
     {(E^L)^{\wedge2g}}
=\frac{(H^{g-1}\cdot\phi^*H)}{(H^g)}
=\frac1{2g}\operatorname{Tr}(\phi\phi').
\]
Choose a symplectic basis for $E^L$ and expand exterior multiplication. The
resulting quadratic form is a sum of squares, hence positive definite.
\end{proof}

\begin{proposition}[Simple endomorphism algebras: initial classification]
\label{mumford-prop-simple-endomorphism-classification}
\dcref{M:ch4:IV.21}
\uses{mumford-thm-rosati-positive,mumford-cor-simple-decomposition}
For simple $X$, $D=\operatorname{End}^0X$ is a division algebra with a
positive involution. If $K$ is its center, $K_0=\{x\in K:x'=x\}$,
$[D:K]=d^2$, and $[K:\mathbb Q]=e$, then $K_0$ is totally real. If
$K\ne K_0$, it is totally imaginary over $K_0$ and the involution is of the
second kind; otherwise it is of the first kind. In the nontrivial Brauer
class case of the first kind, $D$ is a quaternion division algebra.
\source{mumford-abelian-varieties:page-0204}
\source{mumford-abelian-varieties:page-0205}
\source{mumford-abelian-varieties:page-0206}
\end{proposition}
\begin{proof}
Write the real algebra of $K_0$ as
$\mathbb R^{r_1}\times\mathbb C^{r_2}$. Positivity of
$\operatorname{Tr}(xx')$ rules out complex embeddings, so $K_0$ is totally
real. If $K\ne K_0$, positivity forces the quadratic extension to be
totally imaginary. For an involution of the first kind, $D$ is identified
with its opposite algebra, so its Brauer class has order one or two. In the
order-two case its rank over $K$ is four. The canonical quaternion
involution is $x^*=\operatorname{Tr}^0_{D/K}x-x$; Skolem--Noether gives
$x'=ax^*a^{-1}$ with $a^*=\epsilon a$. Extending to the real factors,
positivity excludes $M_2(\mathbb R)$ and leaves Hamilton quaternion factors
with the corresponding positive involutions.
\end{proof}
